% §4.4.4 Case 3c: Vendor disagrees with mainnet applicability / scope characterization
% Swap this file in at Week 23 if vendor disputes the applicability or characterization.
% Probability: 12% (v8.A-final outcome distribution)
% Key handling: structural-not-normative claim in §4.5; reframe as "industry-wide
% open question" linking to §9.

\subsubsection{Coordinated Disclosure and Vendor Response}
\label{sec:pefo-disclosure-3c}

We initiated coordinated disclosure with \DEPLOYEDAPP{} on $D_3$
(\S\ref{sec:disclosure}).

\paragraph{Vendor Response.}
The \DEPLOYEDAPP{} team responded to the disclosure within [X] days. The team
disputes [SPECIFIC ASPECT OF OUR CLAIM — fill from actual response]. Specifically:

``[VERBATIM QUOTE OF VENDOR DISAGREEMENT — fill from actual response].''

The team's position is that [SUMMARY OF VENDOR'S ARCHITECTURAL ARGUMENT —
fill from actual response].

\paragraph{Our Reply.}
We respond to the vendor's disagreement on the following technical grounds.

First, our mainnet applicability claim in \S\ref{sec:pefo-mainnet} is a
\emph{structural} observation: we identify the presence or absence of the three
PEFO conditions (\S\ref{sec:pefo-structural}) in the publicly available contract
code. This is a descriptive, not normative, claim. Whether the conditions
constitute a ``vulnerability'' in the vendor's operational context is a separate
question that we do not adjudicate here.

Second, the specific architectural facts cited in \S\ref{sec:pefo-mainnet}
are verifiable from the on-chain contract state at address
[CONTRACT ADDRESS — fill from actual analysis], block [BLOCK NUMBER]. The
vendor's disagreement concerns [SCOPE / CHARACTERIZATION — fill from actual
response], not the architectural observation itself.

Third, the question of whether cross-rollup sequencing commitments provide
economic atomicity (A4) is an active area of research and design. We present
our analysis as a contribution to this discussion, not as a final determination.
\S\ref{sec:discussion} further explores the design implications.

\paragraph{Documentation.}
The full disclosure exchange is archived in Appendix~\ref{app:disclosure}.
The vendor's response is reproduced [in full / in summary with permission —
choose at camera-ready]. All communications are independently timestamped
via OpenTimestamps proofs.

\paragraph{Disclosure Timeline.}

\begin{table}[h]
\centering
\caption{Disclosure timeline (\DEPLOYEDAPP{}, Case 3c).}
\label{tab:disclosure-timeline-3c}
\begin{tabular}{ll}
\toprule
Event & Date \\
\midrule
Initial outreach & [DATE] \\
Formal disclosure ($D_3$) & [DATE] \\
Vendor response (disagreement) & [DATE] \\
Our reply & [DATE] \\
$D_3 + 90$ day deadline & [DATE] \\
Public disclosure & [DATE] \\
\bottomrule
\end{tabular}
\end{table}
